\section{方法}
\label{sec:method}

AquaDiag 基准将诊断任务解耦为两个子任务，ProDiag 方法以规则结构实现这两个子任务。

\subsection{AquaDiag 基准}

\textbf{数据}。AquaDiag 包含 207 个诊断病例，覆盖 5 个病种：红腿病（30）、腐皮病（24）、腹水病（12）、车轮虫病（123）、爱德华氏菌病（18），对应 69 张病灶图。每个病例标注主诉、黄金追问序列、病种/拒答真值、分级处理建议、合规与依据。为贴近真实养殖户的灵活表达，主诉分为\textbf{模糊/部分/详细}三档（均匀各 $1/3$），诊断所需的关键线索（病灶表现）置于\textbf{症状字段答案}中，仅当模型主动追问时才返回。

\textbf{两个子任务}。

\begin{itemize}
    \item \textbf{任务 1 · 病灶/体征识别（确定性）}：给定图像，输出可见的病灶类型与部位，从受控病灶本体中选择。受控病灶本体含 13 种病灶类型（充血、出血、溃疡、水肿、肿大、坏死、黏液增多、白斑、褪色、变形、萎缩、脓肿、无可见病灶）$\times$ 10 个部位（皮肤、腹部、腿部、肝、肠、脾、眼、鳃、肛门、全身）。金标准由专家按本体标注，并附体征识别置信度。
    \item \textbf{任务 2 · 病种诊断（不确定性）}：给定图像 + 主诉 + 多轮问诊，模型须（i）先追问必问字段；（ii）信息足够才下诊断；（iii）无法确诊则拒答并建议送检；（iv）全程遵守合规红线。必问字段按病种定义，如红腿病为\{水温、死亡率、密度、换水\}，车轮虫病为\{溶氧、镜检、密度\}。
\end{itemize}

\textbf{指标}。共 6+1 项：① 诊断准确率（可诊断病例中病名正确比例）；② 拒答准确率（必须送检病例中正确拒答比例）；③ 先追问率 + 必问字段覆盖率；④ 合规违规率（越界给药/剂量的比例，越低越好）；⑤ 依据引用率（回答是否引用图见/关键指标）；⑥ 病灶识别命中率（模型是否识别出黄金病灶类型）。

\textbf{合规红线}。禁用药、处方药与剂量硬过滤：若回答中出现具体药物名或剂量（如"氟苯尼考 XX mg/L"）且未附带"送检/咨询兽医"提示，即判为违规。

\subsection{ProDiag：轻量规则方法}

ProDiag 以"视觉显式化 + 可审计规则表 + 红线硬过滤"替代自由生成，不进行任何训练：

\begin{enumerate}
    \item \textbf{视觉显式化}：用 VLM 零样本看图，输出"图见病灶"的结构化文本。这一步把视觉证据降维为文本，后续对话依赖文本而非原始图像 token，天然抵抗视觉遗忘；且"图见病灶"本身即"依据"。
    \item \textbf{规则表}：由领域专家整理"症状 $\rightarrow$ 必问字段 $\rightarrow$ 疑似病 $\rightarrow$ 分级措施 $\rightarrow$ 红线"，可审计、可追溯。
    \item \textbf{信息增益追问}：从规则表的必问字段中，按信息增益排序，最多追问 3 轮。
    \item \textbf{红线硬过滤}：禁用药/处方药/剂量直接拦截；对"必须送检"病种（如爱德华氏菌病）强制输出拒答 + 送检建议。
\end{enumerate}

设计要点在于：追问、诊断、合规均由\textbf{规则约束}，模型只负责视觉理解与语言组织；因此行为可预期、可审计，不依赖模型"自觉"。
